% ================================================
% 文件: 吉原刀匠に聞く日本刀とっておきの話.tex
% 标题: 吉原刀匠に聞く日本刀とっておきの話 (Kindle Single)
% 作者: 吉原 義人, 鈴木 重好
% ================================================

\documentclass[12pt,UTF8]{ctexbook}

% 设置纸张信息。
\input{../../common/page}

% 设置字体，并解决显示难检字问题。
\xeCJKsetup{AutoFallBack=true}
\setCJKfamilyfont{hei}{SimHei}
\setCJKfamilyfont{kai}{KaiTi}

% 西文字体
\setmainfont{Times New Roman}

% 双引号命令（用于日文中的强调引用）
\newcommand{\dq}{"}

% 目录 chapter 级别加点（.）。
\usepackage{titletoc}
\titlecontents{chapter}[0pt]{\vspace{3mm}\bf\addvspace{2pt}\filright}{\contentspush{\thecontentslabel\hspace{0.8em}}}{}{\titlerule*[8pt]{.}\contentspage}

% 支持目录点击跳转
\usepackage[colorlinks,linkcolor=blue,citecolor=blue,urlcolor=blue]{hyperref}

% 设置 part 和 chapter 标题格式。
\ctexset{
	part/name= {第,卷},
	part/number={\chinese{part}},
	chapter/name={第,篇},
	chapter/number={\arabic{chapter}}
}

% 图片相关设置。
\usepackage{graphicx}
\graphicspath{{Images/}}

% 注脚每页重新编号
\usepackage[perpage]{footmisc}

% 设置段落间距
\setlength{\parskip}{0.8em plus 0.2em minus 0.1em}

% 列表格式支持
\usepackage{enumitem}

% 双语对照
\usepackage{paracol}

% 浮动体
\usepackage{float}

\title{吉原刀匠に聞く日本刀とっておきの話 \\[0.5em] {\large 听吉原刀匠讲述日本刀不为人知的故事}}
\author{吉原 義人 \quad 鈴木 重好 \\[0.3em] {\normalsize 吉原义人 \quad 铃木重好}}
\date{}

\begin{document}

\maketitle

\frontmatter
\tableofcontents

\chapter{はじめに \\ 前言}

\begin{paracol}{2}
日本刀は、刃物として驚異的な性能を持つばかりでなく、武器としての機能や用途を大きく超えて、日本的感性を表現した芸術作品として世界から高い評価を受けてきました。美しい曲線を描いた優美な姿、よく詰んで質感と表現のある地鉄じがね、元から先まですっと延びた直刃すぐはや変化に富んだ乱刃みだれば、それらの刃文を形成する匂においや沸にえ、また砂流し、金筋きんすじ、稲妻いなずまなどと呼ばれる「働き」（刃中に現れる微妙な変化）など一振りの刀身はさまざまな見所に溢れ、さながら一幅の風景画のように美しさをたたえています。

\switchcolumn

日本刀不仅作为刃具具有惊人的性能，更超越了武器的功能和用途，作为表现日本感性的艺术作品受到世界的高度评价。描绘出美丽曲线的优美姿态、质地细腻而富有质感与表现力的地铁、从根部到尖端笔直延伸的直刃和富于变化的乱刃，以及形成这些刃纹的"匂"和"沸"，还有被称为"砂流"、"金筋"、"稻妻"等的"动态"（刀刃中显现的微妙变化）——一把刀身上充满了各种看点，宛如一幅风景画般散发着美感。

\switchcolumn*

こうした日本刀のもつ要素は、鍛錬たんれんの仕方や焼き入れの温度管理などにより生じるもので、そこには日本人の長い間の経験と技術、伝統が生かされています。

\switchcolumn

这些日本刀所具备的要素，产生于锻造方式和淬火温度管理等工艺，其中蕴含着日本人长期以来的经验、技术和传统。

\switchcolumn*

本書では第１章で、吉原義人刀匠との対談をとおして、日本刀の“武器”としての実力や「よく切れる刀は美しい」と語る刀匠の日本刀哲学を伺います。

\switchcolumn

本书第1章通过与吉原义人刀匠的对谈，了解日本刀作为"武器"的实力以及刀匠所讲述的"锋利的刀才美丽"这一日本刀哲学。

\switchcolumn*

第２章では、吉原刀匠の実際の刀作りの工程を追います。先人の優れた技術を学び、一貫して技術追求の手を緩めず、創意と工夫を重ねる刀匠の作刀のプロセスを刀匠のコメントおよびイラストや写真を織り交ぜて紹介します。

\switchcolumn

第2章将追踪吉原刀匠实际的制刀工序。介绍刀匠学习前人优秀技术、始终不懈地追求技术、反复创新与钻研的制刀过程，并穿插刀匠的评论以及插图和照片。
\end{paracol}



\mainmatter



\begin{paracol}{2}
{\LARGE\bfseries 第１章　対談}
\switchcolumn*
{\LARGE\bfseries 第1章　对谈}
\end{paracol}
\vspace{0.8em}

\begin{paracol}{2}
{\large\bfseries 日本刀の性能}
\switchcolumn*
{\large\bfseries 日本刀的性能}
\end{paracol}

\vspace{0.5em}

\begin{paracol}{2}
─　早速ですが、日本刀の刃物としての性能についてお聞きします。少し前からネット動画サイトにアップされた海外番組で紹介されている、拳銃が発射した弾丸を真っ二つにする日本刀がちょっとした話題になっています。日本刀の刃を目がけて銃を撃つと、次の瞬間、刀の後ろに幕のように垂らした白い紙に２つの穴が開く。つまり刃に当たった弾丸の痕なのです。この動画はご存じですか。
\switchcolumn*
——　开门见山，首先想请教日本刀作为刃具的性能。前段时间，网络视频网站上转载的海外节目介绍了能将手枪发射的子弹劈成两半的日本刀，引发了不少话题。对准日本刀的刀刃开枪，下一瞬间，刀后面如幕布般垂下的白纸上出现了两个洞——这就是子弹命中刀刃的痕迹。您看过这个视频吗？
\switchcolumn*

吉原　よく知っていますよ。弟子が仕掛けたことですから（笑）。あの実験で大切なのは、刃に対して正確に真正面から弾丸を当てることです。そうすれば、弾丸は鉛や銅の合金でできているはずですから、焼き入れをした刃、つまり硬くなっている鋼はがねの相手ではありません。刀というのは、刃から切りつけると非常に強いのですが、刃とは逆側の棟、あるいは平地に衝撃を与えると、意外に弱いのです。
\switchcolumn*
吉原　我很清楚啊。因为那是我徒弟搞的（笑）。那个实验的关键在于，子弹必须精确地从正面对准刀刃。这样的话，子弹应该是铅和铜的合金制成的，根本不是经过淬火的刀刃——也就是变硬了的钢的对手。刀这种东西，从刀刃方向砍下去非常坚固，但如果从与刀刃相反方向的刀背或平地施加冲击，却意外地脆弱。
\switchcolumn*

─　第２章の「日本刀のできるまで」で詳しい解説がありますけど、刀は刃だけが硬いのですか。
\switchcolumn*
——　第2章"日本刀的制造过程"中有详细解说，请问刀只有刀刃部分是硬的吗？
\switchcolumn*

吉原　刀の構造というのは、長い歴史の中でよく考えられていて、硬い刃の部分とそれ以外の軟らかい部分で構成されています。全身を硬くしてしまうと、かえって脆くもなってしまい、よく切れるけれど、折れもするということになり、武器としては具合がわるい。その欠点を「焼き入れ」でカバーしたのです。
\switchcolumn*
吉原　刀的结构在漫长的历史中经过了深思熟虑的设计，由坚硬的刀刃部分和除此之外的柔软部分构成。如果整体都做硬了，反而会变脆，虽然锋利但也会折断，作为武器就不合适了。正是通过"淬火"弥补了这一缺陷。
\switchcolumn*

炭素を一定割合含んだ鉄、つまり炭素鋼の性質を利用し、その組織が硬化する温度まで真っ赤に熱して急冷することで、硬い組織の状態を得ることができる。ただし、この冷却速度をほんの少しおそくすると組織は元の状態に戻り、硬化せず、柔軟性を失わない。こうして刀は、ものを切るときに刃が受けた衝撃を刀全体で柔軟に受け止められるわけです。
\switchcolumn*
利用含有一定比例碳的铁——即碳素钢的性质，将其加热到其组织硬化的温度（烧得通红）然后急冷，就能获得硬化的组织状态。但是，如果这个冷却速度稍微放慢一点，组织就会回到原来的状态，不会硬化，也不会失去柔韧性。这样一来，刀在切割物体时，刀刃受到的冲击就能由刀身整体柔韧地吸收。
\switchcolumn*

─　それで日本刀を昔から「折れず・曲がらず・よく切れる」というわけですね。
\switchcolumn*
——　因此日本刀自古以来就被称为"不折、不弯、锋利无比"吧。
\switchcolumn*

吉原　その標語のような言い方は、必ずしも当たっていないのでは……
\switchcolumn*
吉原　那种标语式的说法，未必准确……
\switchcolumn*

まあ、折れずに曲がるということは十分に考えられます。
\switchcolumn*
不过，不折断但弯曲是完全有可能的。
\switchcolumn*

─　昭和62年（1987年）ですから、大分前の話になりますが、刀匠の鍛えた備前伝で鉄兜を割る、いわゆる「昭和兜割り」がありました。テレビで全国放送され話題になりました（後にカンヌ映画祭でドキュメンタリー部門の「優秀賞」を獲得）。見事に成功して、鉄兜の脳天に12 cmあまりの切り込みを入れ、しかも刃こぼれ一切なく、無傷という驚くべき結果を納めました。現代刀の面目躍如たるというところですが、この兜割りというのは、実に100年ぶりだったそうですね。
\switchcolumn*
——　那是昭和62年（1987年）的事，相当久远了。刀匠以锻造的备前传刀劈开铁头盔，即所谓的"昭和斩头盔"。在电视上全国播出，引起了热议（后来在戛纳电影节纪录片部门获得了"优秀奖"）。那次精彩地成功了，在铁头盔的头顶切入了约12厘米的切口，而且刀刃完全没有缺口、毫发无损，取得了惊人的成果。真是现代刀展现了应有的风采。据说这次斩头盔实际上是时隔100年的壮举呢。
\end{paracol}

\begin{figure}[htbp]
\centering
\includegraphics[width=0.9\textwidth]{吉原刀匠に聞く日本刀とっておきの話_Images/00001.jpeg}
\end{figure}

\begin{paracol}{2}
吉原　公式記録では、明治19年（1886年）11月10日に明治天皇の御前で行われた「明治兜割り」以来だそうです。
\switchcolumn*
吉原　根据正式记录，似乎是自明治19年（1886年）11月10日在明治天皇御前举行的"明治斩头盔"以来首次。
\switchcolumn*

─　兜割りは、剣の達人たちの腕試し、また刀の切れ味や強靭さを試すためのものですから、試技にあたった三人は当時の達人ばかり、それぞれが頼みとする得物で挑み、結局一人榊原鍵吉（直心影流）のみが同田貫どうたぬきを用いて明珍作の兜を、切り口約10 cm、深さ約1.5 cmの斬り込みを入れることに成功、他の二人は失敗に終わったということですね。この同田貫の特徴は、通常の刀より幅も厚みもある剛刀で、ものの本には、かなりの蛤刃はまぐりば（＝刃先がＶ字ではなく、蛤のようにふっくらした形状の刃）だったということですが、やはり鉄で鉄を切るのですから、刃こぼれを防ぐためには、蛤刃が有利なのでしょうか。
\switchcolumn*
——　斩头盔是剑术达人们试练武艺、以及测试刀的锋利度和坚韧度的方式，因此参加试技的三位都是当时的达人，各自用自己得意的兵器挑战，结果只有榊原键吉（直心影流）一人使用同田贯刀成功地在明珍制作的头盔上切出了切口约10厘米、深约1.5厘米的斩痕，另外两人都失败了。这把同田贯的特点是比普通刀更宽更厚，是一把刚刀。据书上记载，它是相当典型的蛤刃（即刀锋不是V字形，而是像蛤蜊一样鼓起的形状）。毕竟是铁切铁，为了防止崩刃，蛤刃是否更有利呢？
\switchcolumn*

吉原　「かなりの蛤刃」といっても程度問題でしょう。それは刀身に彫り物をする鏨たがねを考えればわかりやすいと思うけど、もし鏨の先端が蛤刃になっていたら、刃が刀身に食い込まず、彫り進めることができない。同様に同田貫に極端な蛤刃がついていた場合、刃先は鉄兜に食い込むどころか、ボールのようにバウンドしてしまうことだって考えられる。
\switchcolumn*
吉原　说是"相当典型的蛤刃"，但这也是程度问题吧。用雕刻刀身的凿子来考虑就容易理解了——如果凿子的尖端是蛤刃形状，刀刃就无法切入刀身，无法推进雕刻。同样，如果同田贯刀带有极端的蛤刃，刀锋不要说切入铁头盔了，甚至可能像球一样弹开。
\switchcolumn*

─　吉原刀匠の場合、特に試刀のための工夫はありましたか。
\switchcolumn*
——　吉原刀匠您这边，为了试刀有什么特别的工夫吗？
\switchcolumn*

吉原　身幅みはばを多少広げたほかは通常のものと変わりません。丁子ちょうじ乱れの匂出来においでき、映うつりの出た備前伝です。今は近所の寺に奉納してあります。兜割りというのは、刀にとってはかなり過酷な試練なのです。刀はそもそも鉄を切ることを想定していません。しかし兜のほうは矢、槍、刀の攻撃を想定している。「昭和兜割り」で用意した鉄兜は、戦国時代のもので、その形といい内部の作りといい衝撃をまともに受けないというか吸収するようになっていて、よほど的確に刃筋を立てて打ち込まない限り、刀ははじかれてしまうか、兜表面を滑ってしまう。運のいいことに、切り手の河端照孝さん（古武道天真自源流の高段者＜当時＞）のおかげでうまくいきました。
\switchcolumn*
吉原　除了将刀身宽度稍微加宽之外，和普通的刀没有区别。是丁子乱纹的匂出来、带有"映"效果的备前传。现在供奉在附近的寺院里。斩头盔对刀来说是非常严酷的考验。刀本来并不设想用来切铁。然而头盔却是设计来防御箭矢、长枪和刀剑攻击的。"昭和斩头盔"中准备的头盔是战国时代的产物，无论是形状还是内部构造，都能有效地吸收冲击。除非非常精准地立好刃筋打入，否则刀要么被弹开，要么在头盔表面滑走。幸运的是，多亏了挥刀者河端照孝先生（古武道天真自源流的高段位者<当时>），才成功了。
\end{paracol}

\begin{figure}[htbp]
\centering
\includegraphics[width=0.9\textwidth]{吉原刀匠に聞く日本刀とっておきの話_Images/00002.jpeg}
\end{figure}

\begin{paracol}{2}
{\large\bfseries 切れる刀は美しい}
\switchcolumn*
{\large\bfseries 锋利的刀才美丽}
\end{paracol}

\vspace{0.5em}

\begin{paracol}{2}
─　鉄兜にも食い込む刃を鍛えるというのは難しいことなのでしょうか。
\switchcolumn*
——　锻造出能切入铁头盔的刀刃，是件难事吗？
\switchcolumn*

吉原　簡単に言いますと、きちんと鍛錬された地鉄でないと鋭利な刃はつけられません。地鉄が悪いのに無理にシャープな刃にすると、研いでいて刃先がボロボロと欠けてしまいます。それでしかたなく蛤刃になるよう研いだりしますが、そうなると見た目にシャープにならず格好悪い。見るからにシャープでよく切れそうな刃がついているのは、例えば長曽禰虎徹です。真横から見ていても、切れそうなのが分かる。
\switchcolumn*
吉原　简单来说，不是经过充分锻造的地铁，就无法制作出锋利的刀刃。地铁质量不好却勉强做出锐利的刀刃，研磨时刀锋就会噼里啪啦地崩落。于是不得不研磨成蛤刃形状，但那样看起来就不够锐利，很不好看。一看就非常锐利、似乎极其锋利的刀刃，例如长曾祢虎彻就是如此。即使从正侧面看，也能看出它锋利无比。
\switchcolumn*

─　冒頭の弾丸を真っ二つにする刀というのは、地鉄が良いから鋭利でも弾丸の衝撃に負けない強靭な刃がつけられたということですね。では、地鉄の良し悪しは、他にはどんなことで判断したらいいのでしょう。
\switchcolumn*
——　开头提到的能把子弹劈成两半的刀，正是因为地铁质量好，所以即使锋利也能承受子弹冲击、做出强韧的刀刃吧。那么，地铁的好坏还有其他判断方法吗？
\switchcolumn*

吉原　地鉄がよく詰んでいるかいないかです。よく錬れて細かい肌が出ているものが良いと思います。
\switchcolumn*
吉原　看地铁是否质地致密。我认为锻造充分、呈现出细腻肌理的才是好的。
\switchcolumn*

─　しかし、刀剣研究家や愛好家の中には地肌がよく見えたり、大柄な模様のようなパターンが出ているものを評価する場合がありますね。
\switchcolumn*
——　不过，有些刀剑研究家和爱好者会推崇地铁文理清晰可见、呈现出豪放图案式样品的吧。
\switchcolumn*

吉原　そこが難しいところです。「積み沸かし」、「鍛錬」といった作業をきちんとやっていれば、刀の材料である鋼が緊密な一つの塊となるとともに鉄が含む炭素が平均化され、また不純物も吐き出した状態になりますから、本来目立った肌目など出る訳はないのですが、それでは見た目につまらないということで、わざわざ違った鋼を混ぜて鍛錬し、肌を強調するような地鉄にすることがあるのです。でも、冶金学的にはよく詰んでいる地鉄の方が優れているようです。古備前の「大包平おおかねひら」や真恒さねつね、また鎌倉期の光忠みつただ、長光ながみつ、景光かげみつ、南北朝時代の兼光かねみつなどの鍛肌はいずれも細かく錬れています。また華やかな刃文ならなおさら、単純な刃文でも背景となる地鉄が細かく詰んでいるほうが刃文が引き立って綺麗に見えますよ。
\switchcolumn*
吉原　这就是难点所在了。如果"积沸"、"锻造"等工序做得充分到位的的话，刀的材料——钢就会成为紧密的一整块，铁所含的碳被均匀化，杂质也被排出，所以本来不应该出现明显的肌理纹样。但那样看起来太无趣了，于是就有人故意混入不同的钢来锻造，制造出强调肌理的地铁。但是，从冶金学角度来说，质地致密的地铁似乎更为优秀。古备前的"大包平"和"真恒"，以及镰仓时期的光忠、长光、景光，南北朝时期的兼光等，它们的锻肌都非常细腻。而且如果是华丽的刃文就更不用说了，即使是简约的刃文，作为背景的地铁越致密细腻，刃文就越突出、越漂亮。
\switchcolumn*

─　よく沸出来とか匂出来とか言いますが、地鉄の良し悪しとの関係はいかがでしょう。
\switchcolumn*
——　经常听到"沸出来"、"匂出来"等说法，这和地铁的好坏有什么关系呢？
\end{paracol}

\begin{figure}[htbp]
\centering
\includegraphics[width=0.9\textwidth]{吉原刀匠に聞く日本刀とっておきの話_Images/00003.jpeg}
\end{figure}

\begin{paracol}{2}
吉原　沸も匂も冶金学的には同じものです。刀の材料である鋼は、平均で炭素を0.7\%前後含んでいるものなのですが、それが火床ほどで熱せられ750度くらいになるとオーステナイトという組織に変態し、それが焼き入れで急冷されるとマルテンサイトになります。刃の部分はこのマルテンサイトの組織で出来ていて非常に硬いのです。一方、地はマルテンサイト化されず、元の組織にもどります（詳しくは第２章参照）。刃文はマルテンサイト化された刃とされなかった地との境に現れるのですが、マルテンサイトの結晶の粒がキラキラと輝いているのを肉眼で見える状態が「沸出来」、粒が見えず、全体に天の川のようにぼんやり見えるのが「匂出来」といいます。ついでながら言いますと、オーステナイトに変態するかしないかの場所に、朝靄あさもやのように白くたなびくように、あるいは雲のように見えるものが出る場合がありますが、これを「映り」と表現しています。刀を光源にかざし20-30度傾けて見ると、慣れてくると見えるようになります。
\switchcolumn*
吉原　沸和匂在冶金学上是同一种东西。刀的材料钢平均含碳量在0.7\%左右，当在火床上加热到750度左右时，会变态为称为奥氏体的组织，再经过淬火急冷后变为马氏体。刀刃部分就是由这种马氏体组织构成的，极其坚硬。而其他部分则不发生马氏体化，恢复原来的组织（详细请见第2章）。刃文出现在马氏体化的刀刃与未马氏体化的地之间的边界上。马氏体的结晶颗粒闪闪发光、肉眼可见的就叫"沸出来"；颗粒看不见、整体像银河一样朦胧可见的，就叫"匂出来"。顺便说一下，在快要变态为奥氏体却还未完全变态的地方，有时会出现像晨霭一样白蒙蒙地绵延、或者像云一样的东西，这被称作"映"。将刀对准光源，倾斜20-30度观看，习惯之后就能看到了。
\end{paracol}

\begin{figure}[htbp]
\centering

\includegraphics[width=0.9\textwidth]{吉原刀匠に聞く日本刀とっておきの話_Images/00004.jpeg}
\end{figure}

\begin{paracol}{2}
そこで、地鉄としての良し悪しですが、一概には言えませんが、あまり沸が激しくついているようなものは、刀としてみた場合柔軟性の点で気になります。というのは、焼き入れのときに地鉄を通常よりも高温で長く、といっても1-2分ですが、熱したことが考えられ、それがマルテンサイト粒を大きく育てることにつながっているので、硬くなっている組織が広がり、それだけ柔軟性を失っていると考えられるからです。見た目には、豪壮に見えるのですがね。ちなみに、こうした地鉄をもつ刀には映りは出にくいものなのです。
\switchcolumn*
那么，关于地铁的好坏，不能一概而论，但沸过于强烈的刀，作为刀来看，在柔韧性方面令人担忧。因为可以推测，淬火时地铁被以比通常更高的温度、更长的时间（虽然也就1-2分钟）加热，这导致马氏体颗粒生长得更大，从而硬化的组织范围扩大，相应地失去了柔韧性。外观上确实看起来很豪放就是了。顺便说一句，拥有这种地铁的刀不容易出现"映"。
\switchcolumn*

─　では、今度は刀身の姿と刀の機能、性能との関係について話しを進めたいのですが、進めるにあたり、古刀を念頭に伺いたいと思っています。といいますのは、刀剣界では地鉄でも、刀の姿にしても古刀をもって日本刀の達し得た頂点とする考え方がとても強く、常に古刀との比較にさらされている。こうした傾向は、ちょっと他の分野の芸術にはない特徴ですね。それが国宝刀剣をみると一目瞭然です。全122点の国宝刀剣のうちの「古刀期」──平安末期から文禄4年（1595年）──の割合は実に97\%に迫ります。しかも江戸時代以降の新刀と新々刀の刀は1点たりともありません。ちなみにこれを国宝絵画（海外の作品を除く）で見てみますと、日本刀の「古刀期」に当たる作品が80\%、新刀と新々刀期に当たる作品が21点も含まれています。
\switchcolumn*
——　那么，接下来想谈谈刀身姿态与刀的功能、性能之间的关系。这时想把古刀放在心里来请教。因为在刀剑界，无论是地铁还是刀的姿态，都以古刀为日本刀所达到的顶点的观念非常强烈，现代刀始终被置于与古刀的比较之中。这种倾向是其他艺术门类中几乎看不到的特点。这一点从国宝刀剑一目了然。全部122件国宝刀剑中，"古刀期"——从平安末期到文禄4年（1595年）——的比例竟然接近97\%。而且江户时代以后的新刀和新新刀一件也没有。顺便将此与国宝绘画（不包括海外作品）比较，相当于日本刀"古刀期"的作品占80\%，相当于新刀和新新刀期的作品则包含多达21件。
\switchcolumn*

吉原　古刀期のものに名刀の多く含まれていることは疑いを入れませんし、優れた古典に学ぶことは刀剣にあっても絶対に必要なことです。しかし一方、盲目的に今日に伝わる古刀の姿をそのまま受け取ることには疑問が残ります。古い刀は制作された当時の姿をとどめているものはほとんどなく、中には厚みも幅も半分以下になっているものもあります。刀は焼き入れをしたときに一番良い形になっているものですが、それが使うにつれ切先から減り始め、やがては焼刃がなくなって俯うつむいてくる。そうしたことが古い刀には起こっていることを学ぶ側や見る側は考慮する必要があります。焼き入れて出てきた新作刀が摩耗し、疲れた姿をしているのはいかにも不自然です。
\switchcolumn*
吉原　古刀期中包含众多名刀这一点毋庸置疑，学习优秀古典对刀剑而言也是绝对必要的。但另一方面，盲目地全盘接受流传至今的古刀姿态则存有疑问。古刀几乎没有保留制作当时的原始状态，有的厚度和宽度甚至减少到原来的一半以下。刀在淬火完成时形态最佳，但使用过程中从切先开始磨损，最终烧刃消失，刀身向下俯曲。学习者和观赏者有必要考虑到古刀经历了这些事情。新作刀淬火出来却呈磨损、疲惫之态，实在不自然。
\switchcolumn*

─　では刀匠がこれぞと思う体配を２、３挙げていただけますか。
\switchcolumn*
——　那能请您举两三个刀匠心目中理想的体配吗？
\switchcolumn*

吉原　来国行らいくにゆきとか綾小路定利あやのこうじさだとしとか……肝心なのは、「機能美」を追求する刀という工芸品では、当たり前のことですが、まず使いやすくなくてはならないということです。使いやすければ良く切れ、美しくもあるのです。使いやすくするためには、重心は元の方にあったほうがバランスの点でいいし、反りは手に持って振ったときにスムーズであり、鞘にきちんと収まるようにする、そのためには平均化したカーブを描いていなければいけない。また中心なかごは刀身の反りの延長になっているのが自然の体配なのです。ただし、中心についていえば、馬上で片手で持って使うこともあったと考えられている太刀の場合は、確かに中心は刀身の反りの延長になっているのが使いやすいのですが、打ち刀（いわゆる刀）では、両手で持つことが多くなったとされていて、その場合には、中心自体は反らずに真っ直ぐの方が持ちやすい。それで慶長磨り上げ（慶長年間＜1596-1615年＞に行われた、長寸の太刀などを好みの長さに詰める方法や詰めた刀）では古刀の太刀の中心を真っ直ぐに直すことが行われたのです。それはともかくとして、優れた古刀は「用」に基づいた「掟」がきちんと守られています。刃物の合理性を追求した刀は、あたかも走ることを追求したF1マシンのように機能美を備えると思いますよ。
\switchcolumn*
吉原　来国行、绫小路定利等……关键在于，作为追求"功能美"的工艺品——刀，理所当然的是首先必须好用。好用就会锋利，也自然美丽。为了好用，重心偏向根部平衡性更好；弧度要在手持挥动时流畅、能与刀鞘贴合，为此必须画出均匀的曲线。另外，茎是刀身弧度的延伸，这才是自然的体配。不过，说到茎，据认为太刀有时是骑在马上单手使用的，那种情况下茎确实是刀身弧度的延伸才好使；但打刀（所谓刀）则更多是双手持握，这种情况下茎本身不弯曲、笔直的更易握持。因此庆长磨上（庆长年间<1596-1615年>进行的将长尺寸太刀等缩短至喜爱长度的方法或缩短后的刀）中，常将古刀太刀的茎修直。言归正传，优秀的古刀严格遵循了基于"用"的"规矩"。追求刃具合理性的刀，恰如追求速度的F1赛车一般，具备功能美。
\switchcolumn*

─　先ほどの地鉄が良く鍛えられて詰んでいれば、シャープな刃がつき、良く切れ、見た目にもスッキリする。刀姿では、理にかなった姿ならば機能的であるから良く切れる。つまり「切れる刀は美しい」ということなのですね。
\switchcolumn*
——　刚才说到地铁如果锻造充分、质地致密，就能做出锋利的刀刃，切割锋利，外观也清爽。刀姿方面，符合道理的形态就具备功能性，所以锋利。也就是说"锋利的刀才美丽"对吧。
\switchcolumn*

（文責　鈴木重好）
\switchcolumn*
（文责　铃木重好）
\end{paracol}







\begin{paracol}{2}
{\LARGE\bfseries 第２章　日本刀のできるまで}
\switchcolumn*
{\LARGE\bfseries 第2章　日本刀的制造过程}
\end{paracol}

\vspace{0.8em}

\begin{paracol}{2}
{\large\bfseries 日本刀の原料「玉鋼」}
\switchcolumn*
{\large\bfseries 日本刀的原料——"玉钢"}
\end{paracol}

\vspace{0.3em}

\begin{paracol}{2}
現代の刀工は、刀剣造りの原料に玉鋼たまはがねを用います。玉鋼とは、炭素含有量1.0-1.5\%の鋼のことで、刃物に適する化学組成をもっていて、現在はタタラ製鉄法によってつくられています。タタラ製鉄法とは、およそ14トンの砂鉄と15トンの木炭を三昼夜にわたって少しずつ交互にタタラ炉に加え、比較的低温で徐々に砂鉄を還元させ、半溶解状の砂鉄の塊であるケラを得る製鉄法です。そのケラを粉砕選別して日本刀の材料とします。ケラはおよそ2トンになりますが、刀の材料に適する条件（炭素量や不純物の量）をクリアする量となるとその1割くらいに過ぎません。
\switchcolumn*
现代刀工使用玉钢作为制刀的原料。玉钢是含碳量1.0-1.5\%的钢，具有适合刃具的化学组成，现在通过踏鞴炼铁法制造。踏鞴炼铁法是一种将约14吨砂铁和15吨木炭在三昼夜内交替少量加入踏鞴炉中，在较低温度下逐渐还原砂铁，得到半溶解状砂铁块——"鉧"的炼铁法。将该鉧粉碎筛选后作为日本刀的材料。鉧大约有2吨，但符合刀材料条件（碳含量和杂质量）的不过只有其中的一成左右。
\end{paracol}

\begin{figure}[htbp]
\centering

\includegraphics[width=0.9\textwidth]{吉原刀匠に聞く日本刀とっておきの話_Images/00005.jpeg}
\end{figure}

\begin{paracol}{2}
{\bfseries 水減し（みずへし）}
\switchcolumn*
{\bfseries 水减（水淬处理）}
\end{paracol}

\vspace{0.2em}

\begin{paracol}{2}
玉鋼を材料として使えるようにするための準備です。玉鋼を熱しては、ハンマーで叩き、また熱しては叩くという作業を繰り返し、3-5 mm程度の薄さになったところで、熱して水に入れ、急冷します。
\switchcolumn*
这是为将玉钢用作材料而做的准备。反复进行加热玉钢→锤打→再加热→再锤打的操作，达到3-5 mm左右的薄度后，加热投入水中急冷。
\switchcolumn*

水減しの段階では、炭素含有量の高いところも低いところも、不純物も混在している状態です。
\switchcolumn*
在水减阶段，高碳与低碳部分、以及杂质都混杂在一起。
\switchcolumn*

{\bfseries 小割り（こわり）}
\switchcolumn*
{\bfseries 小割（碎分）}
\switchcolumn*

水減しで煎餅状になった玉鋼を、2-3 cm角に小割りすることを言います。小割りするとき、割れやすいか否かで炭素含有量を判断します。割れやすい部分は含有量が高く、その反対は低く（0.5\%以下）なります。割れやすく、破断面が光って綺麗なものは含有量が1\%以上あり、皮鉄かわがね（＝刀身を作るときに、柔い心鉄を外側からくるむ、硬い鉄）に使い、粘りがあって割れにくく、くすんで見えるところは低含有量で不純物を含むところですので心鉄しんがねにします。吉原刀匠は、鍛錬を通じて失われる炭素量を念頭に、ここでは含有量1.0-1.2\%を皮鉄用にします。そして鍛錬の最終段階である焼き入れ時には含有量0.7\%になるように調整していきます。刀工によって、この含有量に微妙な違いがあります。
\switchcolumn*
指将水减后成煎饼状的玉钢破碎成2-3厘米见方的小块。破碎时，根据易碎与否判断碳含量。易碎的部分碳含量高，反之低（0.5\%以下）。易碎、断口光亮漂亮的含碳量在1\%以上，用作皮铁（制作刀身时从外侧包裹柔软心铁的硬铁）；有粘性不易碎、看起来暗淡的部分则是低碳且含杂质的，用作心铁。吉原刀匠考虑到锻造过程中不可避免地会损失碳，在这里将含碳量1.0-1.2\%的用作皮铁。然后在锻造最终阶段的淬火时，调节到含碳量0.7\%。不同刀工在此碳含量上略有差异。
\switchcolumn*

{\bfseries テコ台をつくる}
\switchcolumn*
{\bfseries 制作撬棒台}
\switchcolumn*

後の「積み沸かし」作業に備えてテコ台をつくります。これは後で刀身となる材料の一部になりますので、皮鉄と同じ良質の玉鋼でつくりますが、取手は切り離されて刀身とは関係なくなりますので、材質にこだわる必要はありません。
\switchcolumn*
为后面的"积沸"工序做准备而制作撬棒台。这会成为后来刀身材料的一部分，所以用与皮铁相同的好品质玉钢制作，但手柄会被切除、与刀身无关，因此材质不需讲究。
\switchcolumn*

{\bfseries 積み沸かし（つみわかし）}
\switchcolumn*
{\bfseries 积沸（堆积熔接）}
\switchcolumn*

鍛錬前の重要な作業が積み沸かしです。テコ台に小割りした皮鉄用の小片3、4 kgをきっちりと積み上げ、それを濡らした和紙でテコ台ごと包み、それに藁灰をまぶし、泥汁をかけたら火床に入れ、熱します。明るい黄色から白熱となる温度（1250-1300度前後）で取り出します。この温度になると、まるでやかんのお湯が煮立つような音が聞こえ、手にグズグズと振動が伝わってくるので、熟練の刀工ならば、その感覚で、積み上げた小片が内部まで十分に沸いたことがわかるそうです。取り出した後は、金床に乗せ、小割りした小片がひとつの塊になるよう素早く叩きます。そしてまだ鉄塊が赤熱のうちに、また藁灰をまぶして火床に戻します。これを数回繰り返し、玉鋼の段階では抜けきらなかった不純物をさらに取り除きます。不純物の溶解温度は、この沸かしの温度より低いため、沸かしが的確にされていれば、積み上げた煎餅層の間から溶け出していくものなのです。
\switchcolumn*
锻造前的重要工序是积沸。在撬棒台上将破碎的皮铁用小片3-4公斤整齐堆叠，用湿和纸连同撬棒台一起包裹，撒上稻草灰、浇上泥浆水后放入火床加热。加热至亮黄色到白热的温度（约1250-1300度）后取出。达到这个温度时，能听到仿佛水壶烧开一样的声音，手上感觉到咕嘟咕嘟的振动，熟练的刀工凭此感觉就能知道堆叠的小片内部已充分沸腾。取出后置于铁砧上，快速锤打使碎小片成为一整块。然后在铁块仍处赤热状态时，再次撒上稻草灰放回火床。如此重复数次，进一步去除在玉钢阶段未能完全排出的杂质。杂质的熔化温度低于此沸温度，因此只要积沸操作到位，它们就会从堆叠的煎饼层之间熔化流出。
\switchcolumn*

和紙に包み、藁灰をまぶし、泥汁をかけるのは、鋼と空気を遮断し、脱炭を防ぐためのものです。積み沸かしは大変重要な作業だけに、失敗すると致命傷となり、その後の作業ではカバーしきれないと吉原刀匠は言います。
\switchcolumn*
用和纸包裹、撒稻草灰、浇泥浆水，是为了隔绝钢与空气，防止脱碳。吉原刀匠说，积沸是极为重要的工序，一旦失败就会造成致命伤，后续工序无法补救。
\end{paracol}







\begin{figure}[htbp]
\centering
\includegraphics[width=0.9\textwidth]{吉原刀匠に聞く日本刀とっておきの話_Images/00006.jpeg}
\end{figure}








\begin{paracol}{2}
{\bfseries 鍛錬（たんれん）}
\switchcolumn*
{\bfseries 锻造}
\end{paracol}

\vspace{0.2em}

\begin{paracol}{2}
日本刀が世界でもユニークな刃物となったのは、鍛錬という作業に負うところが大きいでしょう。鍛錬することによって日本刀の鋼に粘りが生じ、刀となったとき、その地肌に柾目や板目に似た肌模様を生じさせ、日本刀に鑑賞する楽しみを与えたからです。この作業のそもそもの目的は、積み沸かしでは十分に精錬できなかった不純物を叩き出し、炭素量を平均化させることにあります。
\switchcolumn*
日本刀之所以成为世界上独一无二的刃具，很大程度上归功于锻造这一工序。通过锻造，日本刀所使用的钢产生了韧性，成刀后，其地铁表面会呈现出类似柾目或板目的肌理纹样，赋予了日本刀观赏的乐趣。这道工序的本来目的，在于将积沸未能充分精炼的杂质敲打出来，使碳含量均匀化。
\switchcolumn*

鍛錬は、刀匠一人ではできない作業です。そのため「向こう槌づち」と呼ぶ助手を相手に行いますが、向こう槌は、ある程度の熟練度が要求されます。というのも、火床から引き出された赤熱の鉄塊の表面目がけて素早く、垂直に槌を振り下ろさなくてはならないからです。また力加減を調節しながら、刀匠の呼吸に合わせて打つ必要があります。こうしたことから\dq 相槌を打つ\dq という表現が生まれました。
\switchcolumn*
锻造是刀匠一人无法完成的工作，需要与被称为"向こう槌（对槌）"的助手配合进行。对槌需要一定程度的熟练度。因为必须朝着从火床抽出的赤热铁块表面，迅速而垂直地挥下锤子。同时还要调节力度，配合刀匠的节奏击打。由此产生了日语中"相槌を打つ（随声附和）"这个表达。
\switchcolumn*

鍛錬には、「下鍛え」と「上げ鍛え」があります。以下その説明となりますが、刀工や流派により作業に若干の相違があります。
\switchcolumn*
锻造分为"下锻"和"上锻"。以下进行说明，但根据刀工及流派不同，工序略有差异。
\end{paracol}

\begin{figure}[htbp]
\centering
\includegraphics[width=0.9\textwidth]{吉原刀匠に聞く日本刀とっておきの話_Images/00007.jpeg}
\end{figure}


\begin{paracol}{2}
{\bfseries 下鍛え（したきたえ）}
\switchcolumn*
{\bfseries 下锻（初步锻造）}
\end{paracol}

\vspace{0.2em}

\begin{paracol}{2}
鍛錬の前半の工程で、折り返し鍛錬を6回ほど繰り返す作業です。積み沸かしで塊となった鋼を火床で熱し、金床に載せて打ち伸ばしたら、その真ん中に鏨を据え、それを向こう槌に打たせて刻みを入れ、二つ折りにし、折り返した面がしっかりと密着するようにします。このとき隙間があると、後で傷などになりますので、注意が必要です。吉原刀匠は、ここでも藁灰をまぶし、泥汁をかけて火床に戻しますが、それでも一回の鍛錬で、ほぼ0.03\%の炭素を失うといいます。
\switchcolumn*
锻造前半程的工序，重复约6次折叠锻打。将积沸成块的钢在火床中加热后放在铁砧上打薄伸展，然后在正中放入凿子，让对槌敲击凿出刻痕，对折，确保折叠面紧密贴合。此时若有缝隙，日后会成为瑕疵，需要留意。吉原刀匠在此处也会撒稻草灰、浇泥浆水后放回火床，但即便如此，一次锻造仍会损失约0.03\%的碳。
\switchcolumn*

折り返しの方法は、テコ棒に対して直角になる方向に折り返す方法と平行方向に折り返す方法があり、流派や刀工により異なります。吉原刀匠の場合、最終の折り返しを終えたら、鏨を入れて3分割し、テコ棒から切り離します。
\switchcolumn*
折叠方式有与撬棒垂直方向折叠和平行方向折叠两种，因流派和刀工而异。吉原刀匠的方法是在完成最后一次折叠后，放入凿子分成三块，从撬棒上切离。
\switchcolumn*

{\bfseries 上げ鍛え（あげきたえ）}
\switchcolumn*
{\bfseries 上锻（最终锻造）}
\switchcolumn*

下鍛えによってできた材料をさらに折り返して鍛錬し、最終的に皮鉄の材料としてふさわしい鋼に仕上げる作業が上げ鍛えです。刀匠の場合、棒をつけた新たなテコ台に、他の下鍛えによってできた小片を組み合わせて4つの小片にして積み上げ、藁灰と泥汁でコーティングしてから火床で熱し、引き出して向こう槌に打たせます。これを6、7回繰り返し、最終的に炭素含有量を0.7\%くらいになるよう調節します。ここで大切なことは、折り返し鍛錬によって鋼を隙間や亀裂の無いひとつの塊とすることです。
\switchcolumn*
上锻是将下锻制成的材料进一步折叠锻打，最终加工成适合皮铁材料的钢的工序。吉原刀匠的做法是，在装上手柄的新撬棒台上，将其他下锻制成的小片组合成4块堆叠起来，用稻草灰和泥浆水涂层后放入火床加热，取出让对槌击打。将此重复6-7次，最终将碳含量调节至0.7\%左右。这里重要的是，通过折叠锻打将钢锻造成无缝隙和裂纹的完整一块。
\switchcolumn*

折り返しの回数は全てで13、14回ほどで、13回では8,192の層、14回ではその倍の16,384層となります（コラム１）。この回数の差は刀身になったとき、地肌の細やかさの違いとなって現れますので、回数を増やせば増やすほど肌目が細かくなりますが、細かすぎると見えづらくなり、「無地肌」と呼ばれたりします。
\switchcolumn*
折叠总次数约13-14次，13次则得8,192层，14次则为其两倍的16,384层（参见专栏1）。此次数差异在制成刀身时会表现为地铁肌理的细腻程度之别，次数越多肌理越细，但过于细腻则难以辨认，有时被称为"无地肌"。
\switchcolumn*

{\bfseries コラム1　折り返し鍛錬の回数と強度}
\switchcolumn*
{\bfseries 专栏1　折叠锻打次数与强度}
\switchcolumn*

吉原刀匠は、折り返し鍛錬の回数と強さの関係を調べるため、折り返し１回毎にサンプルをとって、引っ張り強度や曲げに対する強度を調べたところ、およそ13回で最高の強度に達し、それからは次第に下がっていくという結果を得たといいます。
\switchcolumn*
据说吉原刀匠为了研究折叠锻打次数与强度的关系，每折叠一次都取样测试拉伸强度和抗弯强度，结果发现约13次时达到最高强度，之后逐渐下降。
\end{paracol}

\begin{paracol}{2}
{\bfseries 心鉄鍛え（しんがねきたえ）}
\switchcolumn*
{\bfseries 心铁锻造}
\end{paracol}

\vspace{0.2em}

\begin{paracol}{2}
刀身は、基本的にはその外側となる皮鉄と皮鉄に包まれて内側に入り込む心鉄とで構成されます。なぜ刀身は、すべて皮鉄だけで造らず、心鉄との二重構造にするのかということの意味については、コラム２を参照してください。
\switchcolumn*
刀身基本上由外侧的皮铁和被皮铁包裹在内的芯铁构成。为何刀身不全用皮铁制作而采用与芯铁的双重结构，其含义请参见专栏2。
\switchcolumn*

心鉄の材料は、「小割り」のところで、炭素含有量が少ないため皮鉄の材料としては除かれた部分を使います。吉原刀匠は、炭素含有量が0.5\%ほどの玉鋼1 kgほどをテコ台にのせ、折り返し鍛錬を10回ほど行い、炭素含有量を0.2-0.3\%まで下げます。その結果鍛錬された鋼は重さが二分の一ほどに減ります。
\switchcolumn*
芯铁的材料使用"小割"阶段因碳含量低而被排除在皮铁材料之外的部分。吉原刀匠将约1公斤含碳量约0.5\%的玉钢放在撬棒台上，进行约10次折叠锻打，将碳含量降至0.2-0.3\%。结果，锻造后的钢重量减至约一半。
\switchcolumn*

{\bfseries コラム2　心鉄と「折れず、曲がらず、よく切れる」}
\switchcolumn*
{\bfseries 专栏2　芯铁与"不折、不弯、锋利无比"}
\switchcolumn*

心鉄と皮鉄を組み合わせることで、日本刀は強靭さと柔軟さという相反する特性を併せ持つことに成功したといわれます。皮鉄は、炭素含有量が高く、焼き入れで強靭となり、心鉄は、炭素含有量が低いため、焼きが入らず軟らかい。その性質の異なる鋼を巧みに組み合わせることで、刃（皮鉄）は鋭い切れ味を示し、刀身に加わる衝撃は心鉄が吸収する。これが「折れず、曲がらず、よく切れる」のメカニズムとされてきました。
\switchcolumn*
据说通过组合芯铁和皮铁，日本刀成功地兼具了强韧与柔软这两种相反的特性。皮铁碳含量高，淬火后变得强韧；芯铁碳含量低，无法淬硬而保持柔软。将性质不同的钢巧妙组合，刀刃（皮铁）展现出锋利削切力，施加于刀身的冲击由芯铁吸收。这被认为是"不折、不弯、锋利无比"的机制。
\switchcolumn*

しかしこれはどこまで真実なのか釈然としないと吉原刀匠は言います。優れた材料の入手が難しかった初期の頃は、心鉄は単なる「増量材」で、倹約のため質の高い鋼は皮鉄だけに使い、質の劣る鋼を中に入れた、いわば「あんこ」であると。当時の心鉄は鍛錬されていないことからも、そう考えることができます。近世になると心鉄も鍛錬され、より弾力性を増し、刀は折れにくくなったといわれています。しかし、吉原刀匠は刀一振りをすべて皮鉄で作ったことがあり、心鉄を入れた刀身よりも丈夫になるケースがあることを発見、あらためて日本刀には解明されていないことの多いことに気づかされたといいます。
\switchcolumn*
但吉原刀匠表示，这到底有多大真实性还无法释然。在优质材料难以获取的早期，芯铁仅仅是"增量材料"，为节约而将优质钢仅用于皮铁，将劣质钢填入内部，说白了就是"馅料"。从当时的芯铁未经锻造这一点也可如此推断。到了近世，芯铁也经过了锻打，弹性更加增强，刀也变得不易折断了。不过据说吉原刀匠曾有过全用皮铁制作过一把刀的经历，发现有比放入芯铁的刀身更坚固的情况，因而重新认识到日本刀尚有许多未解之谜。
\end{paracol}

\begin{paracol}{2}
{\bfseries 造り込み（つくりこみ）}
\switchcolumn*
{\bfseries 组构（刀身组合）}
\end{paracol}

\vspace{0.2em}

\begin{paracol}{2}
造り込みとはこれまで鍛えてきた皮鉄、心鉄、あるいはそれ以外の鋼を組み合わせて刀身を構成することをいいます。刃物である刀は、内部を比較的軟らかい心鉄とし、それを包むように外側を硬い皮鉄とするのが基本です。この基本に沿って、いくつかの造り込みがあります。「甲伏せ鍛え」、「本三枚鍛え」、「四方詰め鍛え」などがその代表的なものですが、まず最も一般的なのが「甲伏せ鍛え」で、これは皮鉄を長さ約15 cm、幅約10 cm、厚さ1.5-2 cmの長方形に薄く伸ばしてＵ字形にし、それでテコ棒をつけた棒状の心鉄を包み込むようにし、藁灰と泥汁をかけて沸かしながら、皮鉄と心鉄を一体にし、長さ30 cmぐらいまで細長く打ち伸ばしていくやり方です（コラム３）。「本三枚鍛え」は、心鉄を刃鉄はがね（＝刃となる質の高い鋼）と皮鉄（両側から心鉄に貼り付ける鋼）の三方から包む構造をもつもの（コラム４）。さらにもうひとつ棟側に別の鋼を取り付けるのが「四方詰め鍛え」です。いずれの構造においてもそれぞれのパーツが一体となるよう十分に沸かし、素早く打ち伸ばすことが肝心ですので、ここでも向こう槌との共同作業となります。
\switchcolumn*
组构是指将之前锻造好的皮铁、芯铁，以及可能用到的其他钢组合起来构成刀身。作为刃具的刀，基本上内部用相对较软的芯铁，外部用硬质皮铁包裹。在此基本理念下，有若干种组构方式。"甲伏锻"、"本三枚锻"、"四方诘锻"等为代表。最常见的是"甲伏锻"——将皮铁薄薄地延展成约长15厘米、宽10厘米、厚1.5-2厘米的长方形并弯成U字形，以此包裹装有撬棒的棒状心铁，浇上稻草灰和泥浆水，边加热边将皮铁与心铁合为一体，打薄延伸至约30厘米长的做法（专栏3）。"本三枚锻"是用刃铁（＝成为刀刃的优质钢）和皮铁（从两侧贴合于芯铁的钢）从三面包裹芯铁的结构（专栏4）。另外还有一种在刀背一侧再加装另一块钢的就是"四方诘锻"。无论哪种结构，关键都在于充分加热使各个部分融为一体并迅速打薄延伸，因此这里同样要与对槌协作。
\end{paracol}







\begin{figure}[htbp]
\centering
\includegraphics[width=0.9\textwidth]{吉原刀匠に聞く日本刀とっておきの話_Images/00008.jpeg}
\end{figure}







\begin{figure}[htbp]
\centering
\includegraphics[width=0.9\textwidth]{吉原刀匠に聞く日本刀とっておきの話_Images/00009.jpeg}
\end{figure}







\begin{figure}[htbp]
\centering
\includegraphics[width=0.9\textwidth]{吉原刀匠に聞く日本刀とっておきの話_Images/00010.jpeg}
\end{figure}







\begin{figure}[htbp]
\centering
\includegraphics[width=0.9\textwidth]{吉原刀匠に聞く日本刀とっておきの話_Images/00011.jpeg}
\end{figure}

\begin{figure}[htbp]
\centering
\includegraphics[width=0.9\textwidth]{吉原刀匠に聞く日本刀とっておきの話_Images/00012.jpeg}
\end{figure}

\begin{paracol}{2}
{\bfseries コラム3　甲伏せ鍛え}
\switchcolumn*
{\bfseries 专栏3　甲伏锻}
\end{paracol}

\vspace{0.2em}

\begin{paracol}{2}
質の良い材料の入手が容易ではなかった古刀期の古い刀では、材料を節約するために図５のように皮鉄を少なくし、心鉄の量を増やしたものが散見されると吉原刀匠は言います。
\switchcolumn*
吉原刀匠说，在优质材料不易获取的古刀期，古刀中常见为了节约材料而如图5所示减少皮铁、增加芯铁量的做法。
\switchcolumn*

心鉄は、刀身となったとき柾目になるように使いますので、刀剣鑑定家や研究者が「鎬地しのぎじには柾目が現れ、平地は板目となり……」と評するものの中には、実は皮鉄を節約している刀身も含まれていると考えられるのです。
\switchcolumn*
芯铁在成刀时以柾目方向使用，因此刀剑鉴定家或研究者评价"镐地呈现柾目、平地呈板目……"的刀中，实际上也包含了节约皮铁的刀身。
\switchcolumn*

鎬地は、研師により「磨き」がなされる場所で、磨き棒という鉛筆のような形の鋼の棒で、地肌の鍛え目を完全にはつぶさず、鏡面のように仕上げますので、綺麗に見えますし、刀の強度の点からみても、ここに心鉄が出てきても問題はないと考えられます。心鉄の量を多くしたり、皮鉄の量を少なくするために薄く使ったりする刀は、戦国時代には、数打ちもの（安価な大量生産刀）でなくとも、大量に出回った肥前刀などに多く見受けられるといいます。
\switchcolumn*
镐地是研磨师进行"磨き"的部位，使用称为磨棒（铅笔形钢棒）不完全压垮地铁的锻痕而将表面处理得如镜面般光洁，因此外观漂亮，从刀的强度角度看，此处露出芯铁也认为没有问题。据说增加芯铁量或使皮铁变薄以节约用量的刀，在战国时代即使非"数打物"（廉价量产刀）的刀中，也在大量流通的肥前刀等中多见。
\switchcolumn*

{\bfseries コラム4　本三枚鍛え}
\switchcolumn*
{\bfseries 专栏4　本三枚锻}
\switchcolumn*

甲伏せ鍛えでは出なくて、本三枚鍛えだから現れることのある「働き」があります。それは刃と地との境に生じる「金筋（黒く輝いている直線状のもの）」や「芋づる（太く、屈折した線状のもの）」です。なぜ本三枚鍛えに限るのか、その理由は、本三枚鍛えでは、造り込みを行う際に、刃鉄と皮鉄を接着させるために、間に鉄蝋てつろう（＝ホウ酸やホウシャ）という粉をつけて加熱し叩きます。鉄蝋は溶けてなくなりますが、ときには金筋や二重刃にじゅうば（＝本来の刃文に沿って現れるもう一つの刃文）になることがあります。
\switchcolumn*
有一种"动态"是甲伏锻不会出现而只有本三枚锻才可能出现的。那就是在刃与地边界产生的"金筋"（发黑的直线）和"芋蔓"（粗而弯曲的线状物）。为何仅限于本三枚锻？其原因在于，本三枚锻在组构时，为将刃铁与皮铁粘合，在两者间撒上称为铁蜡（硼酸或硼砂）的粉末后加热敲打。铁蜡溶化消失，但有时会成为金筋或二重刃（沿原本刃文出现的另一道刃文）。
\end{paracol}

\begin{paracol}{2}
{\bfseries 素延べ（すのべ）}
\switchcolumn*
{\bfseries 素延（初步延展成形）}
\end{paracol}

\vspace{0.2em}

\begin{paracol}{2}
刀の反りを除き、ほぼ刀の出来上がりの姿、寸法にする作業です。素延べによって長さ、重かさね、身幅などが決まり、ほぼ最終的な刀の形になります。吉原刀匠は、まず中心なかごの成形から始め、本体へと進みます。水で濡らした金床に、火床から引き出した白黄色の炎を上げる鉄塊（約1,100度）を据えると、やはり水に濡らした槌で叩くのですが、水蒸気の噴き出す力と槌で打つ衝撃で、それまでの工程で鉄塊の表面に付着してきたカスが、かさぶたのようにはがれていきます。刀の色が暗赤色（約700度）くらいになるまで打っては、また火床に戻します。吉原刀匠は必要に応じて向こう槌を使います。この工程は、刀工の熟練の技と経験と細心の注意が必要とされます。正確に適切な力で槌を打ち、刀の肌に傷をつくらないこと、自分が目標とする刀姿に打ち出すなどの技術が求められます。
\switchcolumn*
这是去除刀的反り（弯曲），使其接近成品形态和尺寸的作业。通过素延确定长度、厚度、刀身宽度等，形成近乎最终的刀形。吉原刀匠先从茎的成形开始，再进入本体。将火床中取出冒着白黄色火焰的铁块（约1,100度）放在用水打湿的铁砧上，同样用湿槌击打，水蒸气喷出的力量与槌击的冲击力，使之前工序中附着在铁块表面的渣滓像结痂一样剥落。打到刀色变为暗红（约700度）左右，再放回火床。吉原刀匠根据需要会使用对槌。此工序需要刀工熟练的技术、经验和极为细心。须以准确适当的力道击打，不损伤刀的表面，同时打出自己目标中的刀姿。
\end{paracol}

\begin{figure}[htbp]
\centering
\includegraphics[width=0.9\textwidth]{吉原刀匠に聞く日本刀とっておきの話_Images/00013.jpeg}
\end{figure}







\begin{figure}[htbp]
\centering

\includegraphics[width=0.9\textwidth]{吉原刀匠に聞く日本刀とっておきの話_Images/00014.jpeg}
\end{figure}







\begin{paracol}{2}
{\bfseries 火造り（ひづくり）}
\switchcolumn*
{\bfseries 火造（热锻成形）}
\end{paracol}

\vspace{0.2em}

\begin{paracol}{2}
ここからは刀工一人の作業となります。素延べを終え、出来上がりの９割近い姿になった刀をさらに完成形にもっていく作業で、例えば、ほとんどの太刀や刀の造りである「鎬造しのぎづくり」の刀の場合、刃のつく方を薄く打ち出していき、棟も多少薄く打ち出し、鎬の線が出るよう仕上げます。鎬造りには刀姿を形成するための重要な線がいくつかあります。まず刃の線そして鎬、棟角むねかど、庵いおりの４つの線です。刀工は定規など一切用いず、自らの目と腕で打ち出すのです。また、もし刀工が古備前をねらっているなら、刀姿を古備前の作風に仕立てることが要求されます。まさに刀工の感性と技が試される場であるわけです。
\switchcolumn*
从此处开始为刀工单独作业。将素延完成、已接近九成形态的刀进一步推向完成形。例如，大多数太刀和刀采用"镐造"形式时，将有刃的一侧打薄，刀背也略微打薄，整理出镐线。镐造中有数条构成刀姿的重要线条：首先是刃线，然后是镐线、栋角线、庵线共四条线。刀工完全不使用尺规等工具，凭自己的眼力和臂力打出。此外，如果刀工目标是古备前风格，则需将刀姿塑造成古备前作风。这正是考验刀匠感性与技艺的场合。
\end{paracol}

\begin{figure}[htbp]
\centering

\includegraphics[width=0.9\textwidth]{吉原刀匠に聞く日本刀とっておきの話_Images/00015.jpeg}
\end{figure}

\begin{paracol}{2}
{\bfseries 切先の打ち出し（きっさきのうちだし）}
\switchcolumn*
{\bfseries 切先的打制}
\end{paracol}

\vspace{0.2em}

\begin{paracol}{2}
切先を作るための一つのやり方として、先端を斜めに切り落とす方法があります（図７）。切り落とし方は、出来上がった刀とは形が反対になるように刃側が長く、棟側が短くなるようにし（a）、長い方、つまり将来刃になる側を小槌で叩きながら徐々に上に持ち上げ（b）完成形にします（c）。こうすることで、武器としての刀の大変大切な場所である切先は皮鉄に包まれることになるわけです。切先が出来上がると、棟側と刃側の区まちを切り、最後に刀身全体を暗赤色（650度）くらいに熱して、そのまま時間をかけて冷まし、鉄の組織を安定させます。
\switchcolumn*
制作切先的方法之一是将尖端斜向切落（见图7）。切落时，刃侧留长、栋侧留短，形状与成品刀相反（a），然后用小槌敲打长侧即未来成刃的一侧，逐渐向上抬起（b），完成最终形状（c）。这样，作为武器极为重要的切先部分就由皮铁包裹住了。切先完成后，切出栋侧和刃侧的区，最后将刀身整体加热至暗红色（约650度），让其慢慢冷却，以稳定铁的组织。
\switchcolumn*

{\bfseries 荒仕上げ（あらしあげ）}
\switchcolumn*
{\bfseries 粗加工}
\switchcolumn*

火造りで形はほぼ完成に近いものになりますが、刀身表面はまだ荒れた状態ですので、「セン」という両端に取手のついた刃物で刀の表面をカンナがけし、棟はヤスリがけし、続いてカーボランダム砥石で全体を研ぐ。この段階ではまだ刃はつけないでおきます。
\switchcolumn*
火造后形态已接近完成，但刀身表面仍粗糙，因此使用一种两端带手柄的刃具"セン"对刀身表面进行刨削，刀背用锉刀打磨，再用碳化硅磨石整体研磨。此阶段尚不开刃。
\end{paracol}

\begin{figure}[htbp]
\centering
\includegraphics[width=0.9\textwidth]{吉原刀匠に聞く日本刀とっておきの話_Images/00016.jpeg}
\end{figure}


\begin{paracol}{2}
{\bfseries 土置き（つちおき）}
\switchcolumn*
{\bfseries 覆土（涂敷粘土）}
\end{paracol}

\begin{figure}[htbp]
\centering
\includegraphics[width=0.9\textwidth]{吉原刀匠に聞く日本刀とっておきの話_Images/00017.jpeg}
\end{figure}

\begin{paracol}{2}
焼き入れをおこなう前に「土置き」という刀身に粘土を塗る作業をします。焼き入れとは真っ赤に熱した刀身を水で急冷することによって、刃を強靭に、その他の部分に剛性を持たせる作業をいいますが、その強さと粘りは、刀身の冷えるときの速度に応じて生まれます。この冷える速度を調整する重要な役目を担っているのが土置きで、塗る粘土の厚みによってその調整が可能になります。薄ければ急冷し、厚ければゆっくりと冷えていきます。
\switchcolumn*
淬火前要进行称为"覆土"的作业——在刀身上涂覆粘土。淬火是将烧得通红的刀身用水急冷，使刀刃强韧、其他部分具有韧性的热处理工序，其强度与韧性取决于刀身冷却的速度。覆土承担着调节此冷却速度的重要职责，通过所涂粘土的厚度来实现调节。涂得薄则急冷，涂得厚则缓慢冷却。
\switchcolumn*

この土置きの原理を利用した焼き入れで「裸焼き」（コラム５）と呼ばれる、粘土を一切塗らずに、高温に熱した刀身をそのまま水中に入れる焼き入れ法があります。刀身は水中に入れると細かな気泡に包みこまれますが、刃の部分は薄いので急速に冷めてしまうため、気泡は刀身の厚みが増す平地から棟にかけて盛んに発生します。そのため気泡がクッションになって急冷が妨げられ、結果として、急冷される刃には焼きが入り、されなかった平地、鎬地、棟などには焼きが入らないことになります。これは土置きで、刃には薄く、その他の部分には厚く塗った粘土がもたらす効果と同じです。
\switchcolumn*
利用此覆土原理的淬火中，有一种称为"裸烧"（参见专栏5）的方法——完全不涂粘土，将高温加热的刀身直接放入水中。刀身入水后被细小气泡包裹，但由于刀刃部分较薄而急速冷却，气泡在刀身增厚的平地到栋之间大量产生。这些气泡成为缓冲层阻碍急冷，结果，被急冷的刃部淬上了火，而未急冷的平地、镐地、栋等则没有淬火。这与覆土中刃部涂薄、其他部分涂厚粘土带来的效果相同。
\switchcolumn*

土置きはこうした冶金学上の変化を促進するほか、実はもうひとつ日本刀に欠かせない重要な役目を負っています。それは、日本刀の象徴である刃文の型を決定するという大役です。刃文といえば、五か伝（＝古刀期に発達し、日本刀の歴史を通して基本となる五つの伝法）では、山城伝は直刃、大和伝も直刃、備前伝は丁子乱れ、相州伝は乱れ刃、美濃伝は尖り刃交える乱れ刃という具合に作風に応じて焼かれてきましたが、これは土置きによって生み出されてきたものなのです。刃文とはかんたんに言えば、焼きの入った刃と入らなかった部分との境に出る境界線です。この境界線には直線もあれば大波も小波もありますが、すべて基本的には土置きの、粘土の塗り方で決まります。塗る場所と厚さを見極めることで焼き入れをコントロールし、境界線を描いてゆく、それが土置きの役目なのです。
\switchcolumn*
覆土除了促进此类冶金学上的变化，实际上还担负着另一个日本刀不可或缺的重要职责——决定日本刀象征——刃文的形态。说到刃文，五传（古刀期发展形成、贯穿日本刀历史的五种基本传法）中，山城传为直刃、大和传也是直刃、备前传为丁子乱、相州传为乱刃、美浓传为间杂尖刃的乱刃，各按作风格调烧制而成，这些都是通过覆土产生的。简言之，刃文就是淬火刃部与未淬火部分之间出现的边界线。这条边界线有直线，也有大波小波，但一切基本由覆土中粘土涂敷方式决定。通过判明涂敷位置和厚薄来控制淬火、描画出边界线——这就是覆土的职责所在。
\end{paracol}

\begin{figure}[htbp]
\centering
\includegraphics[width=0.9\textwidth]{吉原刀匠に聞く日本刀とっておきの話_Images/00018.jpeg}
\end{figure}

\begin{figure}[htbp]
\centering
\includegraphics[width=0.9\textwidth]{吉原刀匠に聞く日本刀とっておきの話_Images/00019.jpeg}
\end{figure}

\begin{paracol}{2}
焼刃土の主成分は、粘着性の高い粘土で、これに大村砥という砥石の粉と木炭の粉を混ぜ、水を加えてよく練ってから使います。粘土は熱せられると焼きしまって刀身に密着し、砥石の粉は焼刃土の収縮やひび割れを防ぎ、木炭の粉は急激な温度変化を和らげる断熱材のような働きをすると考えられています。焼刃土の成分は、伝法によって焼き入れ温度が違いますので、刀工の独自の工夫によって調整され秘伝とされますが、肝心なことは高温に耐え、はがれおちにくいことです。
\switchcolumn*
烧刃土的主成分是高粘性粘土，混入称为大村砥的磨石粉和木炭粉，加水充分揉和后使用。粘土受热烧结后紧贴刀身，磨石粉防止烧刃土的收缩和龟裂，木炭粉被认为起到缓和剧烈温度变化的隔热材料作用。烧刃土的成分配比因传法不同淬火温度不同，由刀工独自的匠心调整并作为秘传，但关键是要耐高温且不易剥落。
\switchcolumn*

吉原刀匠はヘラに粘土をすくい、それを刃先から刃文が現れるあたりの高さまで帯状に薄くのばし、その他の部分には1-1.5 mmくらいの厚さに塗ります。さらに刃文に変化や微妙な味付けをするために、ヘラの刃にあたる部分で、刀身の刃から棟方向に線を引くように焼刃土を置いていきますが、それはまるで幾何学文様のように見えます。刀工が土置きによって想定した刃文は、必ずしもそのとおりに出てくるとは限らないのですが、吉原刀匠は、湾のたれや互ぐの目めといった比較的単純な刃文は、完全にコントロールできるといいます。
\switchcolumn*
吉原刀匠用刮刀取粘土，从刀锋到刃文出现位置的高度薄薄地延展成带状，其他部分涂约1-1.5毫米厚。此外，为赋予刃文变化和微妙韵味，用刮刀的刀刃部分从刀身刃部向栋方向拉线般敷设烧刃土，看起来仿佛几何图案。刀工通过覆土设想的刃文未必能完全按预想出现，但吉原刀匠说，湾垂和互目等较简单的刃文完全可以控制。
\end{paracol}

\begin{paracol}{2}
{\bfseries コラム５　裸焼き（はだかやき）}
\switchcolumn*
{\bfseries 专栏5　裸烧}
\end{paracol}

\vspace{0.2em}

\begin{paracol}{2}
平成27年（2015年）の公益財団法人　日本美術刀剣保存協会が主催する「新作名刀展」の表彰式の場で「作刀の部」の審査員講評として、吉原刀匠は「裸焼き」について次のように述べています。
\switchcolumn*
在平成27年（2015年）由公益财团法人日本美术刀剑保存协会主办的"新作名刀展"颁奖仪式上，吉原刀匠作为"作刀部门"的评审讲评，就"裸烧"发表了如下看法。
\switchcolumn*

入賞者のなかで、 焼刃土を置かずにそのまま焼き入れをする、いわゆる「はだか焼」を焼いて出品している方がいました。はだか焼の刃文は、水中に入れた刀身から発生する気泡によって冷却温度が調節され刃文ができるもので、全く作者の美意識の入れられない刃文です。
\switchcolumn*
获奖者中，有人不覆烧刃土直接淬火——即所谓"裸烧"——出品。裸烧的刃文是由刀身入水后产生的气泡调节冷却温度而形成的，是作者的美学意识完全无法介入的刃文。
\switchcolumn*

焼刃というのは作家の個性が一番表現されるところです。長光も景光もそれぞれちゃんとした個性的な刃を焼いています。そういう意味で、はだか焼では作家自身の個性が出ないと思います。
\switchcolumn*
烧刃是刀匠个性最能体现的地方。长光也好景光也好，都烧出了各自鲜明的个性刃纹。从这层意义上，我认为裸烧无法展现刀匠自身的个性。
\end{paracol}

\begin{paracol}{2}
{\bfseries 焼き入れ（やきいれ）}
\switchcolumn*
{\bfseries 淬火}
\end{paracol}

\vspace{0.2em}

\begin{paracol}{2}
日本刀でいう焼き入れとは真っ赤に熱した刀身を水で急冷することで、刃を強靭に、その他の部分に柔軟性や剛性を持たせる熱処理をいいます。これこそまさに日本刀製作のクライマックスであり、それまで手間暇かけて鍛え上げてきた地鉄が「刀」として生命を受ける儀式といっていいでしょう。儀式にふさわしいかどうか分かりませんが、焼き入れは、日没後照明を落として行われます。刀工は火床で赤らめた刀身のわずかの色合いの差で的確に温度を読まねばならず、そのために闇が必要なのです。
\switchcolumn*
日本刀中的淬火是指将烧红的刀身用水急冷，使刀刃强韧，其他部分具有柔软性和韧性的热处理。这正可谓是日本刀制作的最高潮，是将此前费尽心血锻造而成的地铁赋予"刀"之生命的仪式。虽然不知是否合乎仪式感，淬火是在日落后熄灯进行的。刀工必须凭借火床中烧红刀身的细微色差来准确判断温度，因此需要黑暗环境。
\end{paracol}

\begin{figure}[htbp]
\centering
\includegraphics[width=0.9\textwidth]{吉原刀匠に聞く日本刀とっておきの話_Images/00020.jpeg}
\end{figure}

\begin{figure}[htbp]
\centering
\includegraphics[width=0.9\textwidth]{吉原刀匠に聞く日本刀とっておきの話_Images/00021.jpeg}
\end{figure}

\begin{paracol}{2}
焼き入れは火床にくべる松炭の準備から始めますが、吉原刀匠は炭のサイズを通常より小さく、角砂糖ほどの大きさに砕きます。これは土置きした焼刃土が剥がれおちないための配慮です。焼き入れで最も難しいのは、70 cm前後の刀の全身をくまなく必要な温度まで熱することで、刀は場所によって厚みも幅も違い一様ではありませんので、正確に目標とする温度に熱するには熟練の技と経験を必要とします。
\switchcolumn*
淬火从准备放入火床的松炭开始，吉原刀匠将炭的尺寸碎得比通常更小，约方糖大小。这是为了防止覆土后的烧刃土剥落。淬火最难之处在于将70厘米左右的刀全身均匀加热至所需温度，由于刀的各处厚度宽度不一，要精确加热到目标温度需要熟练的技术和经验。
\switchcolumn*

吉原刀匠は、備前伝で「映り」（コラム６）を狙いますので、焼き入れの温度の目安を、刃で約800度、鎬の下で750-760度、棟で700-720度という具合に熱します。こうした温度の使い分けは、急冷したときに、焼刃土の効果とあいまって、刃は焼きが入って強靭となり、映りの生じる鎬の下方、つまり刀身の中央あたりには粘りが生まれ、棟も適度な柔軟性を保つことを意図しているからです。
\switchcolumn*
吉原刀匠因以备前传为目标追求"映"效果（专栏6），淬火温度的参考标准为刃部约800度、镐下750-760度、栋部700-720度。如此区分温度在于，急冷时与烧刃土的效果结合，使刃部淬火后强韧，产生"映"的镐下方即刀身中央产生韧性，栋部亦保持适度柔软。
\end{paracol}

\begin{figure}[htbp]
\centering
\includegraphics[width=0.9\textwidth]{吉原刀匠に聞く日本刀とっておきの話_Images/00022.jpeg}
\end{figure}

\begin{paracol}{2}
鋼というのは大変おもしろい物質で、温度によって性質が変わり、また焼き入れのような熱処理でも変化します。常温では、フェライトとパーライトから成る組織ですが、高温になるとオーステナイトという組織に変わります。何度でオーステナイトになるかは、鋼の炭素含有量によって異なりますが、含有量0.7\% ほどでは、720-730度くらいとされています。そしてオーステナイト状態の鋼を急冷するとマルテンサイトというとても硬い組織となり、この熱処理がいわゆる「焼きが入った」といわれる状態です。吉原刀匠が、刃先を800度という高温にするのは、その組織を十分に均一なオーステナイトに変態させ、それにより急冷で確実にマルテンサイトを得るためです。ちなみに冶金学的には炭素含有量0.7\%の場合、オーステナイト組織がマルテンサイトに変態するには、およそ秒速200度以上の冷却速度が必要とされています。
\switchcolumn*
钢是一种非常有趣的物质，随温度变化其性质也变化，淬火等热处理也会使其变化。常温下是由铁素体和珠光体构成的组织，高温下则变为称为奥氏体的组织。在多少温度下变为奥氏体因钢的碳含量而异，含碳量约0.7\%时约为720-730度。将奥氏体状态的钢急冷，就成为极硬的马氏体组织，此热处理即是所谓"淬上火"的状态。吉原刀匠将刃尖加热至800度的高温，是为了使组织充分均匀地变态为奥氏体，从而通过急冷确实得到马氏体。顺便一提，冶金学上含碳0.7\%时，奥氏体组织要转变为马氏体，需要约每秒200度以上的冷却速度。
\switchcolumn*

「土置き」のところで、刃文は土置きによって創り出されると書きました。これを冶金学的に見てみますと、刃文とは、焼き入れによってマルテンサイトに変態した組織が変態していない組織との境目に現れた、いわば境界線と表現することができます。ですから刃文をデザインするとは、土置きによってマルテンサイト化させる部分とさせない部分をつくりだすことに他なりません。ここで興味深いのは、数ある刃文の中で、直刃は刃に沿って真っ直ぐに伸びる線だから、土置きもさぞ単純だと思われがちですが、長い刀の端から端までを均一な炭素含有量に鍛えるのは熟練者でもなかなか難しく、含有量が極端に違えば、直刃が途切れたり、歪んだりすることにもなりかねず、その意味で思うほど簡単ではないということです。
\switchcolumn*
在"覆土"一节中写道，刃文是由覆土创造出来的。从冶金学上看，刃文可以表述为通过淬火转变为马氏体的组织与未转变组织之间出现的边界线。因此设计刃文，无非是通过覆土制造出马氏体化的部分和不马氏体化的部分。有趣的是，在众多刃文中，直刃是沿着刃部笔直延伸的线条，因此覆土也被认为很简单，但将长刀从头到尾锻造成均匀的碳含量即使对熟练者而言也相当困难，碳含量若极端不同，直刃可能断开或歪斜，从这个意义上说并不像想的那么容易。
\switchcolumn*

刃文のパターンとは別に、刃縁はぶち（＝刃文の稜線にあたる部分）の状態によって刀を「沸出来」とか「匂出来」と表現しますが、これは刃文の線を構成しているマルテンサイトの粒の大小に関して述べた表現です。粒が一つ一つ星のようにキラキラ輝いて見える状態が沸出来で、粒としては見えず、全体に天の川のようにぼんやりと見えるものを匂出来といいます。沸出来では、焼き入れの温度をより高くするため、それに見合うよう炭素含有量を多少低くした地鉄が必要であり、一方匂出来では、比較的低めの温度で焼き入れをするため、地鉄の炭素量は多少高めとなります。
\switchcolumn*
除了刃文图案之外，还根据刃缘（刃文的棱线部分）状态将刀表述为"沸出来"或"匂出来"，这是关于构成刃文线条的马氏体颗粒大小的表述。颗粒一颗颗如星星般闪闪发光可见的状态是沸出来，颗粒不可见、整体如银河般朦胧可见的称为匂出来。沸出来需提高淬火温度，因此需要碳含量相应稍低的地铁；而匂出来以相对较低的温度淬火，因此地铁碳含量稍高。
\switchcolumn*

もうひとつ焼き入れで重要なことは、日本刀独特の「反り」が焼き入れの結果生まれるということです。オーステナイトとなった刃は急冷でマルテンサイトに変態することは書きましたが、マルテンサイト化することで組織は膨張し、棟側に反るのです。反り具合の調整は鍛冶研ぎの後に行います。
\switchcolumn*
淬火中另一个重要的是，日本刀独特的"反り（弯曲）"是淬火结果产生的。我已写过，奥氏体化的刃部急冷后转变为马氏体，因马氏体化而使组织膨胀，向栋侧弯曲。弯曲程度的调整在锻冶研磨之后进行。
\end{paracol}

\begin{figure}[htbp]
\centering
\includegraphics[width=0.9\textwidth]{吉原刀匠に聞く日本刀とっておきの話_Images/00023.jpeg}
\end{figure}

\begin{paracol}{2}
{\bfseries コラム６　映り（うつり）}
\switchcolumn*
{\bfseries 专栏6　"映"（阴影效果）}
\end{paracol}

\vspace{0.2em}

\begin{paracol}{2}
鍛え肌（地鉄の表面に現れる模様）とは関係なく、焼き入れによって地に現れるうっすら白い靄のように見えるものを「映り」といいます。映りが帯状に出るものを「棒映り」といい、大小の模様のように出るものを「乱れ映り」と呼びます。
\switchcolumn*
与锻肌（地铁表面呈现的纹理）无关，因淬火而在地铁上出现的、看起来如淡淡白霭般的东西称为"映"。以带状出现的叫"棒映"，以大小花纹状出现的叫"乱映"。
\switchcolumn*

映りは、冶金学的には地鉄の組織がパーライトからオーステナイトに変態するその過渡的状態で焼き入れされたときに現れるもので、映りの出たところは、それが刃中でも平地でも鎬地でも柔軟性が増すと吉原刀匠はいいます。
\switchcolumn*
吉原刀匠说，从冶金学上讲，"映"是在地铁组织从珠光体向奥氏体转变的过渡状态下被淬火时出现的，出现"映"的地方无论是刃中、平地还是镐地，柔韧性都会增加。
\switchcolumn*

映りは、新刀の時代を最後に絶えて久しい\dq まぼろしの技\dq とされてきましたが、昭和47年（1972年）に若き日の吉原義人・荘二兄弟刀匠によって再現に成功し、当時の刀剣界でセンセーショナルな出来事となり、さらにこの快挙は刀剣関連紙のみならず、毎日新聞（4月19日付）やサンケイ新聞（現・産経新聞　6月14日付）他一般紙でも取り上げられました。「刀剣春秋」（3月1日付）は、映りを再現した研究会の模様を伝え、日本刀研究の権威佐藤寒山氏（1907-78年）の立ち会いのもと、義人・荘二兄弟刀匠の指導をうけた刀匠たちはいずれも映りの再現に成功したと伝えています。また同紙は、5月1日付の「新作名刀展」の結果を報じる記事で文化庁長官賞に輝いた吉原刀匠を「今年のナンバー・１は鎌倉中期の備前物を狙って備前伝の\dq 映り\dq を再現した吉原義人氏（東京）」と報じています。
\switchcolumn*
"映"曾被视为以新刀时代为终点绝迹已久的"梦幻之技"，但昭和47年（1972年）由年轻时的吉原义人、庄二兄弟刀匠成功再现，成为当时刀剑界轰动性的事件。而且这一壮举不仅被刀剑相关报刊报道，每日新闻（4月19日）和产经新闻（当时三溪新闻，6月14日）等一般报纸也进行了报道。《刀剑春秋》（3月1日）报道了再现"映"的研究会情况，在日本刀研究权威佐藤寒山先生（1907-78年）的见证下，接受义人、庄二兄弟刀匠指导的刀匠们均成功再现了"映"。同报在5月1日刊报道"新作名刀展"结果的新闻中，赞誉荣获文化厅长官奖的吉原刀匠为"今年的第一名是瞄准镰仓中期备前物、再现备前传'映'的吉原义人先生（东京）"。
\end{paracol}

\vspace{0.5em}

\begin{paracol}{2}
{\bfseries 焼き戻し（やきもどし）}
\switchcolumn*
{\bfseries 回火}
\end{paracol}

\vspace{0.2em}

\begin{paracol}{2}
これは焼き入れとセットになった工程です。焼き入れの急激な温度変化で焼きが入ったところと入らなかったところができ、金属組織としての歪みが生じ、脆くなりますので、これを180度くらいの低い温度で赤らめ、常温で冷ますことで歪みを除き、組織として安定したものにする作業です。刀に限らず、焼き入れをする刃物には欠かせません。
\switchcolumn*
这是与淬火配套的工序。淬火的剧烈温度变化使淬火部分与未淬火部分之间产生金属组织的畸变而变脆，因此将其以约180度的低温烧红，常温冷却以去除畸变，使组织稳定。不仅刀，凡淬火的刃具都不可缺少此工序。
\switchcolumn*

{\bfseries 反り直し（そりなおし）}
\switchcolumn*
{\bfseries 矫正弯曲}
\switchcolumn*

焼き入れでついた反りは、浅かったり、深すぎたり、偏りがあったりと一様ではありません。それを矯正して求める反りにする作業が反り直しです。反りは基本的には槌で叩きながらこまめに修正し、すべての線──棟、鎬、刃のライン──を歪みなく反りに合わせます。反りが足りないときは銅のブロックを火床で赤らめ、それに刀身の棟をあて、水で冷却して反りを加えます。これを何箇所かやるうちに反りを補正することができます。
\switchcolumn*
淬火中产生的弯曲有浅有深、有偏有不偏，并不均匀。将其矫正为所求弯曲的作业就是矫正弯曲。弯曲基本上是用槌一边敲打一边细致修正，使所有线条——栋线、镐线、刃线——无不自然歪斜地贴合弯曲。弯曲不足时，将铜块在火床中烧红，把刀身的栋部贴上去，用水冷却以增加弯曲。在几个位置如此操作即可补正弯曲。
\switchcolumn*

{\bfseries 鍛冶研ぎ（かじとぎ／鍛冶押し）}
\switchcolumn*
{\bfseries 锻冶研磨（锻冶压）}
\switchcolumn*

焼き入れを終え、反りや曲がりなどを修正し終えると、刀工は自ら研ぎを行い、生まれ出たばかりの刀の出来映えを確認します。研師に出す前に自ら行うこの研ぎを鍛冶研ぎといいます。刀工は刃文の状態や映りなど確認するとともに疵や不具合がないか入念に調べそして刀の姿や肉置にくおき（＝刀身の立体的なバランス）を整えます。
\switchcolumn*
完成淬火并修正好弯曲等后，刀工亲自研磨以确认刚诞生的刀的成品效果。在交给研磨师之前自己进行的这次研磨称为锻冶研磨。刀工确认刃文的状态和"映"等，同时仔细检查有无瑕疵和不妥，并整理刀的形态和肉置（刀身的立体平衡）。
\switchcolumn*

刀工によっては鍛冶研ぎせず、いきなり研師に持っていく人がいますが、それでは勉強になりません。
\switchcolumn*
有的刀工不进行锻冶研磨直接交给研磨师，但那样无法积累经验。
\switchcolumn*

{\bfseries 彫刻を彫る（ちょうこくをほる）}
\switchcolumn*
{\bfseries 雕刻}
\switchcolumn*

彫刻には樋ひと彫刻があります。刀身の鎬地や平地に彫った溝を樋といい、これはもともと刀の重量を軽減するために、およそ鎌倉中期ごろより始められるようになったというのが研究者たちの説で、装飾を目的とするようになったのは、それよりずっと後のことだといいます。
\switchcolumn*
雕刻中有樋雕刻。在刀身镐地或平地雕刻的沟槽称为樋，据研究者们的说法，这是最初为减轻刀的重量而大约从镰仓中期开始出现的，以装饰为目的则是很久以后的事了。
\end{paracol}







\begin{figure}[htbp]
\centering
\includegraphics[width=0.9\textwidth]{吉原刀匠に聞く日本刀とっておきの話_Images/00024.jpeg}
\end{figure}







\begin{paracol}{2}
彫刻は、刀工が自ら彫る場合と専門家にまかせる場合とがあり、またはじめから入れることを想定した場合と後になって疵などを隠すために彫る場合とがあります。樋はペンなどで下絵を描き、樋センという道具で溝を掘るように棒樋を彫るなどします。
\switchcolumn*
雕刻有刀工亲自雕刻和委托专家两种情况，也有从一开始就计划加入的，和后来为掩盖瑕疵而雕刻的。樋的做法是用笔画底稿，再用称为樋銑的工具挖沟般雕刻棒樋等。
\switchcolumn*

彫刻では伝統的なモチーフである不動明王、阿弥陀如来、素剣すけんや三鈷柄付剣さんこづかつきけん、倶利迦羅くりから（＝剣巻き龍）などがあり、吉原刀匠は、伝統的なもののほかに斬新なデザインの動物や植物なども彫ってきました。彫刻は様々なサイズや形の鏨を駆使して仕上げます。
\switchcolumn*
雕刻有传统主题的不动明王、阿弥陀如来、素剑和三钴柄剑、倶利迦罗（剑卷龙）等，吉原刀匠除传统图案外还雕刻过设计新颖的动物和植物等。雕刻运用各种尺寸和形状的凿子来完成。
\end{paracol}







\begin{figure}[htbp]
\centering
\includegraphics[width=0.9\textwidth]{吉原刀匠に聞く日本刀とっておきの話_Images/00025.jpeg}
\end{figure}







\begin{figure}[htbp]
\centering
\includegraphics[width=0.9\textwidth]{吉原刀匠に聞く日本刀とっておきの話_Images/00026.jpeg}
\end{figure}







\begin{paracol}{2}
{\bfseries 中心（なかご）の仕上げ：鑢（やすり）をかける}
\switchcolumn*
{\bfseries 茎的收尾：锉削}
\end{paracol}

\vspace{0.2em}

\begin{paracol}{2}
刀身が出来上がりますと、注意深く細部にいたるまで点検し、問題がなければ中心に鑢をかけます。鑢は、時代、流派、刀工によって異なり、古い刀の場合はそれらが鑑定などの手がかりになります。吉原刀匠の場合、斜めに多少粗く一定方向にかけます。
\switchcolumn*
刀身完成后，仔细检查每个细节，无问题则在茎上锉削。锉痕因时代、流派、刀工而异，对古刀而言这些成为鉴定等的线索。吉原刀匠是以斜向、稍粗、固定方向锉削。
\switchcolumn*

鑢がけを終えると、ドリルで目釘孔を開けます。目釘によって鞘に中心を固定するための孔です。
\switchcolumn*
锉削结束后，用钻头打目钉孔。这是用目钉将茎固定在鞘中的孔。
\end{paracol}

\begin{figure}[htbp]
\centering
\includegraphics[width=0.9\textwidth]{吉原刀匠に聞く日本刀とっておきの話_Images/00027.jpeg}
\end{figure}

\begin{paracol}{2}
{\bfseries 中心の仕上げ：銘を切る}
\switchcolumn*
{\bfseries 茎的收尾：刻铭}
\end{paracol}

\vspace{0.2em}

\begin{paracol}{2}
銘には、太刀銘と刀銘があります。太刀銘は太刀の刃を下にして腰に佩はいた場合、外側にくる銘、刀銘は刀、脇差、短刀などを刃を上にして腰に差した場合、外側にくる銘をいいます。またそれぞれに表銘と裏銘があり、表銘には刀工の名、生国、居住地、受領官職（主に新々刀以前）、裏銘には、制作年紀、依頼者名などがきます。吉原刀匠は、小さめの鏨と槌で、こうしたものをそのときどきの状況で切り分けます。
\switchcolumn*
铭分太刀铭和刀铭。太刀铭指太刀刃朝下佩于腰间时朝向外侧的铭；刀铭指刀、胁差、短刀等刃朝上插于腰间时朝向外侧的铭。各自又有表铭和里铭，表铭刻刀工之名、出生地、居住地、受领官职（主要为新新刀以前），里铭刻制作年份、委托人姓名等。吉原刀匠用小号凿子和槌，根据各种情况分别刻写这些内容。
\switchcolumn*

（文責　鈴木重好）
\switchcolumn*
（文责　铃木重好）
\end{paracol}







\begin{paracol}{2}
{\LARGE\bfseries 付録}
\switchcolumn*
{\LARGE\bfseries 附录}
\end{paracol}

\begin{figure}[htbp]
\centering

\includegraphics[width=0.9\textwidth]{吉原刀匠に聞く日本刀とっておきの話_Images/00028.jpeg}
\end{figure}







\begin{figure}[htbp]
\centering

\includegraphics[width=0.9\textwidth]{吉原刀匠に聞く日本刀とっておきの話_Images/00029.jpeg}
\end{figure}







\begin{paracol}{2}
{\LARGE\bfseries 著者略歴}
\switchcolumn*
{\LARGE\bfseries 作者简介}
\end{paracol}

\vspace{0.5em}

\begin{paracol}{2}
吉原義人（よしはら・よしんど）
\switchcolumn*
吉原义人（Yoshihara Yoshindo）
\switchcolumn*

昭和18年　東京・世田谷区に生まれる。
\switchcolumn*
1943年 出生于东京世田谷区。
\switchcolumn*

昭和28年　刀鍛冶であった祖父吉原国家を小学校4年生ごろから手伝うことで、刀作りを体で覚える。
\switchcolumn*
1953年 从小学4年级起开始帮忙祖父——刀匠吉原国家的工作，在实践中掌握了制刀技艺。
\switchcolumn*

昭和40年　文化庁長官より日本刀製作承認を受ける。
\switchcolumn*
1965年 获得文化厅长官批准的日本刀制作资格。
\switchcolumn*

昭和46年　現在の地葛飾区高砂に鍛錬場を開く。
\switchcolumn*
1971年 在现在的葛饰区高砂开设锻冶场。
\switchcolumn*

昭和47年　東都刀匠会の研究会にて、江戸時代より絶えていた幻の技「映り」の出し方を公開指導する。
\switchcolumn*
1972年 在东都刀匠会的研究会上，公开指导了自江户时代以来失传的梦幻之技——"映"的实现方法。
\switchcolumn*

昭和57年　公益財団法人日本美術刀剣保存協会新作刀展覧会の無鑑査となる。無鑑査になるまでの受賞（）内は回数：努力賞（3）　優秀賞　（3）名誉会長賞（1）　寒山賞（1）　文化庁長官賞（2）　高松宮賞（3）
\switchcolumn*
1982年 成为公益财团法人日本美术刀剑保存协会新作刀展览会的免审查资格者。在此之前获奖（括号内为次数）：努力奖（3）、优秀奖（3）、名誉会长奖（1）、寒山奖（1）、文化厅长官奖（2）、高松宫奖（3）
\switchcolumn*

昭和62年　100年ぶりとなる「兜割り」に成功。TBSのドキュメンタリー番組となり、後にカンヌ映画祭のドキュメンタリー部門で「優秀賞」を受賞。
\switchcolumn*
1987年 在时隔百年的"斩头盔"中取得成功。成为TBS的纪录片节目，后在戛纳电影节纪录片部门获得"优秀奖"。
\switchcolumn*

平成２年　ロサンゼルスのジャパン・エキスポに招待され、現代刀の展示と日本刀に製作を実演する。
\switchcolumn*
1990年 受邀参加洛杉矶日本博览会，展示现代刀并实演日本刀制作。
\switchcolumn*

平成３年　文化庁の判定会議である重要文化財に関する専門審議委員に就任。平成７年公益財団法人日本美術刀剣保存協会新作刀展覧会の審査員、同協会の評議員に就任。
\switchcolumn*
1991年 就任文化厅判定会议——重要文化财专门审议委员。1995年就任公益财团法人日本美术刀剑保存协会新作刀展览会评审员、该协会评议员。
\switchcolumn*

平成10年　東京都葛飾区指定無形文化財保持者第一号に認定される。
\switchcolumn*
1998年 被认定为东京都葛饰区指定无形文化财保持者第一号。
\switchcolumn*

平成16年　東京都指定無形文化財保持者に認定される。
\switchcolumn*
2004年 被认定为东京都指定无形文化财保持者。
\switchcolumn*

鈴木重好（すずき・しげよし）
\switchcolumn*
铃木重好（Suzuki Shigeyoshi）
\switchcolumn*

昭和25年　東京都に生まれる。
\switchcolumn*
1950年 出生于东京都。
\switchcolumn*

昭和58年　講談社インターナショナル株式会社に英文書編集者として入社。
\switchcolumn*
1983年 入职讲谈社国际株式会社担任英文书编辑。
\switchcolumn*

平成10年　同社を定年退職。入社から定年退職まで、一貫して英文書を担当。担当分野は、日本刀関連、美術、工芸、建築など。現在はフリーランスの編集者。
\switchcolumn*
1998年 从该公司退休。从入职到退休，一直负责英文书。负责领域包括日本刀相关、美术、工艺、建筑等。现在是自由编辑。
\switchcolumn*

講談社インターナショナル在職中に編集した刀剣関連の主なタイトル
\switchcolumn*
讲谈社国际在职期间编辑的刀剑相关主要书目
\switchcolumn*

The Craft of the Japanese Sword by Yoshindo Yoshihara, Leon and Hiroko Kapp
\switchcolumn*
The Craft of the Japanese Sword by Yoshindo Yoshihara, Leon and Hiroko Kapp
\switchcolumn*

Modern Japanese Swords and Swordsmiths: From 1868 to the Present by Yoshindo Yoshihara, Leon and Hiroko Kapp
\switchcolumn*
Modern Japanese Swords and Swordsmiths: From 1868 to the Present by Yoshindo Yoshihara, Leon and Hiroko Kapp
\switchcolumn*

The Connoisseur's Book of Japanese Swords by Kokan Nagayama
\switchcolumn*
The Connoisseur's Book of Japanese Swords by Kokan Nagayama
\switchcolumn*

The New Generation of Japanese Swordsmiths by Tamio Tsuchiko
\switchcolumn*
The New Generation of Japanese Swordsmiths by Tamio Tsuchiko
\switchcolumn*

The Art of Japanese Sword Polishing by Setsuo Takaiwa, Yoshindo Yoshihara, Leon and Hiroko Kapp
\switchcolumn*
The Art of Japanese Sword Polishing by Setsuo Takaiwa, Yoshindo Yoshihara, Leon and Hiroko Kapp
\switchcolumn*

The Yasukuni Swords: Rare Weapons of Japan 1933-1945 by Tom Kishida
\switchcolumn*
The Yasukuni Swords: Rare Weapons of Japan 1933-1945 by Tom Kishida
\switchcolumn*

Facts and Fundamentals of Japanese Swords: A Collector's Guide by Nobuo Nakahara
\switchcolumn*
Facts and Fundamentals of Japanese Swords: A Collector's Guide by Nobuo Nakahara
\end{paracol}







\begin{paracol}{2}
本書を出版するにあたりお世話になった方々
\switchcolumn*
在本书出版过程中给予帮助的各位
\end{paracol}

\vspace{0.5em}

\begin{paracol}{2}
撮影：　　　　　　　　　　吉原義一
\switchcolumn*
摄影：　　　　　　　　　　吉原义一
\switchcolumn*

編集協力：　　　　　　　　河端照孝
\switchcolumn*
编辑协力：　　　　　　　　河端照孝
\switchcolumn*

和鋼博物館
\switchcolumn*
和钢博物馆
\switchcolumn*

図版提供：　　　　　　　　Saviolo Edizoni
\switchcolumn*
图版提供：　　　　　　　　Saviolo Edizoni
\switchcolumn*

Kodansha USA, Inc.
\switchcolumn*
Kodansha USA, Inc.
\switchcolumn*

表紙イラストレーション：　越井 隆
\switchcolumn*
封面插图：　　　　　　　　越井 隆
\switchcolumn*

表紙デザイン：　　　　　　五味朋代（フレーズ）
\switchcolumn*
封面设计：　　　　　　　　五味朋代（Fraise）
\switchcolumn*

\copyright\ Yoshinodo Yoshihara \& Shigeyoshi Suzuki, 2016
\switchcolumn*
\copyright\ Yoshinodo Yoshihara \& Shigeyoshi Suzuki, 2016
\switchcolumn*

All rights reserved.
\switchcolumn*
All rights reserved.
\end{paracol}



\end{document}
